% ==================== Chapter 7: Conclusion ====================
\newpage\section{Conclusion and Future Work} [Estimated: 1,500 words]

\subsection{Summary of Contributions}
This paper presents a GDPR compliance warning framework based on Android permission auditing. The core contributions are:

\begin{itemize}
\item A non-intrusive, low-power permission call frequency audit using \texttt{AppOpsManager.getHistoricalOps()}, WorkManager, and Room, achieving <150 MB memory and <2\% extra battery drain.
\item An intelligent decision engine with ``expert baseline + dynamic threshold'' that adapts to application categories and historical variance, outperforming fixed-threshold approaches.
\item A modular, auditable compliance engine that explicitly supports GDPR Articles 5, 7, and 9, producing verifiable audit trails, supporting regulatory framework switching, and enabling independent verification of consent revocation effectiveness.
\item An automated violation simulator that enables systematic, randomised stress testing of the detection logic, providing statistical confidence in precision and recall metrics.
\item A practical validation methodology combining Monte Carlo logical injection, automated simulation-based pressure testing, and a single real-device smoke test, balancing academic rigour with engineering feasibility.
\end{itemize}

\subsection{Key Findings}
Experimental results show precision $\ge 99\%$ and recall $\ge 98\%$ in 1000 Monte Carlo rounds. The 50-cycle automated simulation confirms consistent detection performance across diverse random configurations. The single 24-hour smoke test successfully captured malicious background calls and triggered warnings. The sensitivity analysis confirms the 3-sigma choice provides optimal balance between precision and recall.

\subsection{Future Work}
\begin{enumerate}
\item \textbf{Privacy policy integration}: incorporate NLP techniques to automatically extract permission usage claims from privacy policies and compare against actual runtime behaviour.
\item \textbf{Cross-platform extension}: extend the framework to emerging operating systems like HarmonyOS NEXT, studying its microkernel-based permission auditing mechanisms.
\item \textbf{Streaming anomaly detection}: introduce real-time detection of sudden permission usage spikes, reducing detection latency below 24 hours.
\item \textbf{Large-scale user study}: conduct a broader user study to evaluate the cognitive impact of the tiered notification strategy and refine the verification-friction balance.
\item \textbf{Federated learning for baselines}: explore privacy-preserving aggregation of baseline statistics across devices to improve threshold accuracy without centralising sensitive data.
\end{enumerate}

\subsection{Concluding Remarks}
This work provides a practically feasible and academically sound technical path for operationalising GDPR's data minimisation principle on mobile devices, with transparency and auditability as first-class design goals. The integration of automated randomisation into the validation methodology sets a new standard for rigorous evaluation of privacy compliance tools.